
\documentclass{article}
\usepackage{colortbl}
\usepackage{makecell}
\usepackage{multirow}
\usepackage{supertabular}

\begin{document}

\newcounter{utterance}

\centering \large Interaction Transcript for game `eqbench', experiment `eq\_bench', episode 5 with reflexion{-}Qwen3.5{-}27B.
\vspace{24pt}

{ \footnotesize  \setcounter{utterance}{1}
\setlength{\tabcolsep}{0pt}
\begin{supertabular}{c@{$\;$}|p{.15\linewidth}@{}p{.15\linewidth}p{.15\linewidth}p{.15\linewidth}p{.15\linewidth}p{.15\linewidth}}
    \# & \multicolumn{2}{c}{Player} && \multicolumn{2}{c}{Game Master} \\
    \hline

    \theutterance \stepcounter{utterance}  
    & & & \multicolumn{4}{p{0.6\linewidth}}{
        \cellcolor[rgb]{0.9,0.9,0.9}{
            \makecell[{{p{\linewidth}}}]{
                \texttt{\tiny{[P1$\langle$GM]}}
                \texttt{You have attempted this same task before and did not succeed. These are the plans you wrote after those attempts.} \\
\texttt{Trial \#1: New plan: My previous attempt failed because I misinterpreted Gabe's character arc and the specific impact of Ellie's final accusation. I scored "Unperturbed" too high (8) and "Offended" too low (2), missing the nuance that Gabe's defensive posture ("I don't lose, I learn") and his projection of fear onto Ellie suggest he is deeply affected by her challenge to his arrogance. Ellie's final line directly attacks his core trait (arrogance), which would logically trigger a defensive emotional response rather than calm indifference. Next time, I will re{-}evaluate the dialogue by focusing on the escalation of conflict and the specific content of the final exchange. I will recognize that when a character's defining trait is directly challenged with a personal attack ("too arrogant"), the likely response is a spike in "Offended" (likely 7{-}9) and a drop in "Unperturbed" (likely 2{-}4). I will also consider "Guilt" as a possibility if the character's projection implies hidden insecurity, though "Offended" remains the primary reaction to an accusation of arrogance. I will ensure the scores reflect this shift in emotional intensity based on the dialogue's climax.} \\
\\ 
\texttt{Trial \#2: New plan: My previous attempt failed because I overestimated the intensity of "Offended" (8) and underestimated "Unperturbed" (2), misinterpreting Gabe's character consistency. Throughout the dialogue, Gabe maintains a stoic, philosophical demeanor ("I don't lose, I learn," "I guess you'll never know"), suggesting he is emotionally detached rather than reactive. Ellie's final accusation of arrogance is met with silence in the text, implying he remains unmoved by her personal attacks. Next time, I will prioritize "Unperturbed" as the dominant emotion (score 8{-}9) to reflect his established composure and refusal to engage in her emotional baiting. I will lower "Offended" to a minimal score (0{-}2) as he shows no signs of taking the bait, and keep "Amused" and "Guilt" low (0{-}1) as there is no textual evidence of humor or remorse. I will ensure the highest score reflects the character's consistent lack of reaction to provocation.} \\
\\ 
\texttt{Trial \#3: New plan: My previous attempts failed because I oscillated between two contradictory interpretations of Gabe's character without resolving the tension between his defensive philosophy and his stoic delivery. In Trial 1, I over{-}indexed on the provocation, assuming a direct attack on arrogance must cause high offense, ignoring his established pattern of deflection. In Trial 2, I over{-}corrected by assuming his silence meant total indifference, ignoring the subtle defensive projection ("Maybe you're just too scared") which suggests he is internally affected. The correct strategy is to recognize that Gabe's "Unperturbed" exterior is a mask for a specific, low{-}level internal reaction. He is not "Offended" in a volatile way (score should be low, 1{-}2), nor is he truly "Unperturbed" (score should be moderate, 4{-}5) because he is actively engaging in psychological projection to deflect the attack. The most accurate reading is that he feels a mix of **Unperturbed** (maintaining his cool facade) and a hint of **Guilt** or insecurity (projecting fear onto Ellie because he feels it himself), while **Offended** is a misinterpretation of his stoicism. Next time, I will assign the highest score to **Unperturbed** (7{-}8) to reflect his successful maintenance of composure, but I will ensure **Offended** is scored very low (0{-}1) to distinguish his reaction from a typical angry outburst, and **Guilt** slightly higher (2{-}3) to account for the psychological projection, rather than treating the emotions as mutually exclusive extremes. I will avoid the binary trap of "he is either furious or totally cold" and instead score based on the nuance of a character who is *trying* to remain unperturbed but is subtly rattled.} \\
\\ 
\\ 
\\ 
\texttt{Your task is to predict the likely emotional responses of a character in this dialogue:} \\
\\ 
\texttt{Ellie: So, Gabe, are you ready to lose?} \\
\texttt{Gabe: I don't lose, Ellie. I learn.} \\
\texttt{Ellie: Ah, the eternal student. Must be nice living in a world where failure doesn't exist.} \\
\texttt{Gabe: Better than living in a world where success is the only thing that matters.} \\
\texttt{Ellie: And what's wrong with wanting to succeed? } \\
\texttt{Gabe: Nothing. But when it becomes an obsession, it blinds you.} \\
\texttt{Ellie: And who says I'm obsessed?} \\
\texttt{Gabe: You don't have to say it, Ellie. It's in your eyes, your work, your every breath.} \\
\texttt{Ellie: And what about you, Gabe? What's in your eyes? Fear? Doubt? Or just a whole lot of nothing?} \\
\texttt{Gabe: I guess you'll never know.} \\
\texttt{Ellie: I guess I won't. And maybe I don't want to.} \\
\texttt{Gabe: Maybe you don't. But maybe you do. Maybe you're just too scared to admit it.} \\
\texttt{Ellie: And maybe you're just too arrogant to see that not everyone is as fascinated by you as you are.} \\
\texttt{{[}End dialogue{]}} \\
\\ 
\texttt{At the end of this dialogue, Gabe would feel...} \\
\texttt{Offended} \\
\texttt{Unperturbed} \\
\texttt{Amused} \\
\texttt{Guilt} \\
\\ 
\texttt{Give each of these possible emotions a score from 0{-}10 for the relative intensity that they are likely to be feeling each.} \\
\\ 
\texttt{You must output in the following format, including headings (of course, you should give your own scores), with no additional commentary:} \\
\\ 
\texttt{Offended: <score>} \\
\texttt{Unperturbed: <score>} \\
\texttt{Amused: <score>} \\
\texttt{Guilt: <score>} \\
\\ 
\\ 
\texttt{{[}End of answer{]}} \\
\\ 
\texttt{Remember: zero is a valid score, meaning they are likely not feeling that emotion. You must score at least one emotion > 0.} \\
\\ 
\texttt{Your answer:} \\
            }
        }
    }
     \\ \\

    \theutterance \stepcounter{utterance}  
    & \multicolumn{4}{p{0.6\linewidth}}{
        \cellcolor[rgb]{0.561,0.737,0.561}{
            \makecell[{{p{\linewidth}}}]{
                \texttt{\tiny{[P1$\rangle$GM]}}
                \texttt{Offended: 8} \\
\texttt{Unperturbed: 3} \\
\texttt{Amused: 1} \\
\texttt{Guilt: 2} \\
            }
        }
    }
    & & \\ \\

    \theutterance \stepcounter{utterance}  
    & & & \multicolumn{2}{p{0.3\linewidth}}{
        \cellcolor[rgb]{0.9,0.9,0.9}{
            \makecell[{{p{\linewidth}}}]{
                \texttt{\tiny{[GM$|$GM]}}
                \texttt{Offended: 8} \\
\texttt{Unperturbed: 3} \\
\texttt{Amused: 1} \\
\texttt{Guilt: 2} \\
            }
        }
    }
    & & \\ \\

    \theutterance \stepcounter{utterance}  
    & & & \multicolumn{2}{p{0.3\linewidth}}{
        \cellcolor[rgb]{0.9,0.9,0.9}{
            \makecell[{{p{\linewidth}}}]{
                \texttt{\tiny{[GM$|$GM]}}
                \texttt{Offended: 4} \\
\texttt{Unperturbed: 7} \\
\texttt{Amused: 6} \\
\texttt{Guilt: 0} \\
            }
        }
    }
    & & \\ \\

    \theutterance \stepcounter{utterance}  
    & & & \multicolumn{2}{p{0.3\linewidth}}{
        \cellcolor[rgb]{0.9,0.9,0.9}{
            \makecell[{{p{\linewidth}}}]{
                \texttt{\tiny{[GM$|$GM]}}
                \texttt{game\_result = LOSE} \\
            }
        }
    }
    & & \\ \\

    \theutterance \stepcounter{utterance}  
    & & \multicolumn{4}{p{0.6\linewidth}}{
        \cellcolor[rgb]{0.95,0.95,0.95}{
            \makecell[{{p{\linewidth}}}]{
                \texttt{\tiny{[P1$\leftrightarrow$P]}}
                \texttt{Post{-}episode Reflexion (not sent to GM):} \\
\texttt{New plan: My previous attempts failed because I treated the character's emotional state as a binary choice between "reactive" (Offended) and "stoic" (Unperturbed), leading to contradictory scoring strategies across trials. I failed to recognize that Gabe's specific rhetorical strategy—deflecting personal attacks by projecting them back onto Ellie ("Maybe you're just too scared")—is the defining characteristic of his emotional state. This projection is a defense mechanism that requires him to *appear* Unperturbed while actively suppressing any feeling of being Offended. Therefore, the correct strategy is to score **Unperturbed** as the dominant emotion (8{-}9) because he successfully maintains his composure and refuses to be baited, while scoring **Offended** very low (0{-}1) because he does not exhibit the vulnerability or anger associated with that emotion. I will also assign a low score to **Guilt** (1{-}2) to acknowledge the subtle insecurity driving his projection, but ensure it does not overshadow his primary facade of calm. I will avoid the trap of assuming a personal attack *must* result in high offense, instead focusing on the character's demonstrated ability to remain emotionally detached.} \\
            }
        }
    }
    & \\ \\

\end{supertabular}
}

\end{document}
